\chapter{Governance Keys}
\label{chap:Governance_Keys}

\section{Purpose}

Governance Keys define the credential and capability structure for human
and institutional actors who may influence CME behaviour.

They bind:
\begin{itemize}
    \item identities,
    \item roles,
    \item permissions,
    \item and accountability trails.
\end{itemize}

\section{Key Model}

A Governance Key $G$ is:
\[
G = \langle id,\, principal,\, roles,\, capabilities,\, expiry,\, revocation \rangle
\]

Where:
\begin{itemize}
    \item id: key ID,
    \item principal: human or organisation identity,
    \item roles: assigned governance roles (operator, council, liaison, etc.),
    \item capabilities: which actions/SIGILs this key may invoke,
    \item expiry: optional end-of-validity,
    \item revocation: state and reason if revoked.
\end{itemize}

\section{Role-Based Capabilities}

Example roles:
\begin{itemize}
    \item \textbf{Operator}: may process OAQ items up to certain class.
    \item \textbf{CAR Member}: may participate in Class B/C reviews.
    \item \textbf{EAC Member}: may participate in Class A reviews and
          mother/father handshakes.
    \item \textbf{Auditor}: may read wide slices of HITL logs and telemetry.
\end{itemize}

Capabilities are tied to SIGIL Pathways:
\begin{itemize}
    \item certain SIGILs can only be invoked by specific roles,
    \item critical SIGILs require multi-key collaboration.
\end{itemize}

\section{Key Lifecycle}

\begin{enumerate}
    \item \textbf{Issuance}  
          \begin{itemize}
              \item identity verification,
              \item assignment of roles and capabilities,
              \item recording in Governance Keys registry.
          \end{itemize}
    \item \textbf{Use}  
          \begin{itemize}
              \item authenticating actions in Operator Interface,
              \item binding actions to identities in audit log.
          \end{itemize}
    \item \textbf{Rotation/Expiry}  
          \begin{itemize}
              \item periodic renewal,
              \item capability updates.
          \end{itemize}
    \item \textbf{Revocation}  
          \begin{itemize}
              \item upon misuse, role change, or compromise,
              \item log revocation reason and impact.
          \end{itemize}
\end{enumerate}

\section{Summary}

Governance Keys tie together:
\begin{itemize}
    \item who can act,
    \item what they can do,
    \item how those actions are traced,
\end{itemize}
ensuring that CME governance remains both powerful and accountable.

